\chapter{Data and complete-trip representation}
\label{ch:native-data}

This chapter establishes the observations used to learn traffic motion, the independent requests used to evaluate synthesis, and the representations presented to the generator and GAF readers. The supplied CNR-associated exports provide the original application context. Native MoST demand supplies the complete simulated trips for the current training experiment. A new, prospective SUMO run supplies reference observations only after the generation requests have been fixed. Keeping these roles separate prevents a simulator reconstruction or a classification result from being mistaken for evidence about generated trajectories.

\section{Original exports and coordinate system}

The project began with two supplied tables. The principal export is \nolinkurl{midimost_with_class2.csv}; the secondary export is \nolinkurl{vehicle_positions_TS_New.csv}. The principal MoST export contains 4,840,789 observations from 2,685 identifiers and six participant classes. The secondary TS export contains 2,189,591 observations from 12,061 identifiers but only one recorded class. They differ in spatial support, trajectory length and class composition. Pooling their rows would consequently change several aspects of the learning population at once. The current study uses the principal export to characterize the original domain and does not treat the secondary table as an additional source of interchangeable labelled trips.

The canonical record is \texttt{time,id,x,y,type}. Time is expressed in seconds and positions use local planar coordinates in metres. The principal table also contains simulator-provided speed and a label named \texttt{class2}; the class is mapped to \texttt{type} for the common interface. For comparisons between observed and generated motion, velocity is calculated from positions and elapsed time on both sides. This avoids comparing a simulator's internal speed field with a different numerical derivative of generated coordinates.

No original SUMO network, route file or configuration accompanied the exports. A geometric check against the public Monaco SUMO Traffic network supported using its unchanged projected coordinate frame \cite{most}. All 4,841 systematically sampled principal positions were within five metres of a lane centreline; only 11.2\% of the corresponding 2,190-point secondary sample met that threshold. This check supports the map alignment of the principal source. It neither recovers its exact simulation settings nor proves that all original observations satisfy the current route decoder. The coordinates are not interpreted as raw longitude and latitude.

\section{Choosing a learning domain that preserves class structure}

An initial expansion generated class-permitted random routes on the same network. This increased the number of simulated observations, but a road permission does not recreate a transport service. In particular, permitting a vehicle to behave as a bus does not specify the public line, stops or activity pattern associated with bus observations in the original domain. A source comparison was therefore conducted before continuing the complete-trip generator.

The comparison used the same six classes and matched class counts: 1,983 training and 257 validation entities in each source. The sources were the supplied CNR observations, random-demand SUMO trajectories, and native MoST demand. The inherited spatial encoding, ResNet18 architecture and fixed training budget were held constant. Two training seeds were used, with the final epoch retained rather than a validation-selected checkpoint. These were classifiers trained within this project; the group did not supply classifier weights. Table~\ref{tab:native-source-diagnosis} summarizes the relevant diagnostic evidence.

\tab{Matched six-class source comparison. Values are validation macro F1 in percent. The GAF row averages two fixed classifier seeds. These are recognition results on observed trajectories, not diffusion-generation scores.}{native-source-diagnosis}{\begin{tabularx}{\linewidth}{Xrrr}\toprule
Readout & CNR & Random demand & Native MoST\\\midrule
Spatial geometry &79.12&34.85&76.95\\
Physical motion &88.08&80.57&90.35\\
Inherited Hilbert--GAF &59.88&30.15&58.69\\\bottomrule
\end{tabularx}}

The spatial classifier recovered from 30.15\% macro F1 on random demand to 58.69\% on native demand, compared with 59.88\% on CNR. The geometry readout changed in the same direction, while the motion readout showed that the random-demand data still contained substantial class-associated movement information. Thus, low spatial GAF performance in that domain could not be assigned to the image encoder alone. The evidence supports a mismatch in the demand domain. It does not isolate the causal contributions of service routes, stops, traffic density and activity chains, and the similar CNR and native averages are not an equivalence test.

The native source was fixed to public MoST revision \texttt{b29b2f65}, with a network file identical to the one already used by the project. Its demand includes 24 public bus services and 181 stops, together with activities, parking and background traffic. This is a documented new use of public demand, not a reconstruction of the unavailable original CNR configuration. The matched source study also shares many demand templates with CNR: 204 of 257 validation entities have the same source identity, and the native validation bus lines are represented in training. Its scope is therefore recovery of domain structure under controlled sources, not external validation on new lines.

The practical decision is to train the present generator on native-demand observations. The six-class 58.69\% result is not assigned to the new five-class feedback reader, whose complete-trip cohort differs. Nor does recovery of recognition imply a benefit from GAF selection. That benefit is assessed only after generation, using the new reference population and independent physical metrics.

Participant labels and simulator vehicle classes are separate metadata. SUMO vehicle types specify physical and behavioural parameters whose defaults vary with class~\cite{sumovtypes}. The native source types are retained rather than replacing their permission classes with label names. In particular, assigning the output label taxi or ride-hailing does not itself change a source passenger vehicle into a taxi-permission vehicle. Differences learned from this source remain simulator- and demand-dependent.

\section{Event-supported complete vehicle trips}

The native training source is a single SUMO realization, identified as pilot 300001, with a $0.25$-second simulation step. Its augmented output records the events needed to distinguish a completed leg from an arbitrary observation segment. For ordinary vehicles, a leg extends from departure or the end of an activity stop to the next activity stop or final arrival. For buses, intermediate service stops remain inside the leg. They are part of the movement to be learned, rather than boundaries at which a bus trip is shortened. The earlier observation study's longest continuous segment and the present event-supported leg consequently refer to different data units.

The event-based extraction retained 1,644 legs across six classes after rejecting 15 coordinate-jump cases above the fixed 100~m/s check. The present vehicle experiment excludes the pedestrian walking stages and keeps all 1,051 accepted vehicle legs. It does not additionally filter them by success of the manually implemented map constructor. A route observed in training can legitimately form a supervised route--motion pair even when constructing a comparable route from public map inputs requires an engineering adapter.

The inherited parent-level partition gives 925 training legs and 126 development legs, with counts shown in Table~\ref{tab:native-vehicle-cohort}. Every retained vehicle leg in this cohort has a distinct parent. Known identities retain their earlier training or development assignment, and source parent identities assigned to the original test partition were excluded using the stored split metadata. The new experiment does not read the frozen test trajectories for fitting, selection or evaluation. Distinct parents nevertheless share one traffic realization and public demand; the 1,051 legs must not be treated as 1,051 independently simulated worlds.

\tab{Complete vehicle-leg cohort for the generator and both GAF readers. The same 925 training and 126 development examples are used for the three learned components.}{native-vehicle-cohort}{\begin{tabular}{lrrr}\toprule
Participant class & Training & Development & Total\\\midrule
Bus &81&12&93\\
Motorcycle &226&43&269\\
Private car &457&50&507\\
Bicycle &92&14&106\\
Taxi or ride-hailing &69&7&76\\\midrule
Total &925&126&1,051\\\bottomrule
\end{tabular}}

Completion is supported by simulation events, while the numerical trajectory ends at its last recorded coordinate. In the six-class extraction, 1,123 legs have a last observation $0.25$ seconds before the ending event and 521 coincide with it. No coordinate is invented between the last observation and the event. Durations in the learning table therefore describe the observed extent of a completed leg. This distinction is retained when comparing generated endpoint arrival with sampled simulator output. The extraction also preserves short legs, long bus services and zero-displacement intervals; it does not discard difficult durations to simplify sequence modelling.

\section{Separating route geometry from physical time}

For a training trajectory with positions $\mathbf{x}_i$, cumulative observed arc length is
\begin{equation}
 a_0=0,\qquad a_i=\sum_{j=1}^{i}\lVert\mathbf{x}_j-\mathbf{x}_{j-1}\rVert_2,
 \qquad s_i=\frac{a_i}{a_n}.
 \label{eq:native-observed-progress}
\end{equation}
For nonzero-length legs, $(s_i,t_i)$ gives the progress target. Repeated positions create zero progress increments and retain their elapsed time. The observed spatial polyline retains its nonduplicate vertices in double precision, so reconstructing the coordinates at their recorded progress values does not require a fixed 2,049-point spatial approximation. A 129-point equal-arc summary is used only to condition the neural network. That summary is not the polyline used for final coordinate interpolation.

The choice of temporal target was informed by a controlled reconstruction comparison on the 1,644-leg six-class cohort. With the same target-derived spatial approximation, 257 relative-time points, one-second controls and quarter-second controls were decoded at the original observation times. The average native-speed reconstruction MAE was respectively 0.2614, 0.1960 and 0.0093~m/s; native-acceleration MAE was 1.0514, 1.1542 and 0.0314~m/s$^2$. Improving one-second summaries did not ensure preservation of subsecond movement. These numbers concern known-target reconstruction, including pedestrians and a common spatial approximation. They are not generated-trajectory errors or a result of the new five-class model.

This comparison motivates a variable-length $0.25$-second target and a matched 257-point relative-time control. The final generation clock includes a fractional last interval when the sampled duration is not divisible by $0.25$. Both models use continuous cumulative progress; they do not normalize disconnected temporal blocks and join them afterwards. Their supervision uses the same route--motion pairs and the same class-conditioned duration prior. Masking and length-aware batches make different sequence lengths practical, and the training loss averages over physical time within a trajectory before averaging trajectories. A long bus service is therefore not assigned extra loss weight solely because it contains more samples.

Retaining a finer target removes one representational limitation, but it also increases the learning problem and computation. Equal optimization steps do not mean equal numbers of processed temporal values. A model can preserve the possibility of subsecond variation and still learn an inaccurate or excessively variable distribution. For this reason the generation comparison reports both one-second motion summaries and native-interval speed and acceleration, rather than using reconstruction fidelity as a substitute for synthesis quality.

\section{Spatial and velocity images from the same observations}

The spatial reader inherits the original project's position-only transformation \cite{original,salvadori,ferreira}. Coordinates are normalized with the frozen spatial bounds, mapped through the existing Hilbert routine, and processed by its sequence squeezing and duplication rules to obtain 224 values. The inherited image pipeline produces GASF, GADF and MTF channels, with frozen channel scaling. The shorthand ``spatial GAF'' in this thesis refers to that complete inherited RGB pipeline; the MTF channel is not itself a Gramian angular field. No explicit velocity, acceleration or elapsed-duration feature is appended to the spatial reader.

Position-only input does not imply that all effects of movement disappear. Recorded position order and the pattern of repeated locations depend on traversal. However, they are filtered by the spatial grid and sequence processing. A whole-dataset class distinction can therefore coexist with weak differences among same-route candidates. The comparison retains the original transformation so that its feedback value is measured under a documented lineage, rather than silently replacing it with a displacement-based or velocity-assisted classifier.

The velocity reader starts from the same timestamps and coordinates. One-second increments of cumulative coordinate arc length provide a speed sequence, summarized at 128 relative-time locations. The existing velocity-image procedure uses a fixed 35~m/s scale and GASF, GADF and MTF channels, followed by the inherited resizing and normalization. It is a separate representation, not the original spatial method. Values exceeding the scale and spatial coordinates exceeding frozen bounds are counted. In particular, 128 samples of one-second speed summaries do not retain every quarter-second acceleration fluctuation, even though the new generator can produce such fluctuations.

Both readers are newly trained ResNet18 networks \cite{resnet}, initialized from ImageNet and fitted to the same five vehicle classes on the same 925 training legs. Each uses the fixed final epoch of a 50-epoch run; the 126 development legs describe recognition and probability quality without choosing a best epoch. The velocity reader in this experiment predicts participant class, not the six speed-mode labels used in an earlier development study. Their matched cohort and labels support a fair representation comparison, while the separate inputs preserve the question of what spatial and velocity information each contributes.

\section{Prospective requests and independent simulator references}

Formal generation uses a different interface from supervised training. Before a new simulation, the experiment freezes 100 requests: one for each of the 24 public bus services and 19 for each of the other four vehicle classes. Non-bus requests use eligible public single-leg plans from development-side parents, with no activity parking inside the requested leg. The fixed request specifies the planned edge sequence, participant type, start and end lane positions, and any public bus stops. These are available conditions, not facts recovered from a realized future trajectory.

The independent reference run uses seed 330201, retains native background demand, and inserts the targets at fixed new departure times between simulation seconds 18,000 and 18,600. Target rerouting is disabled to match the fixed-plan task, while background traffic remains active. SUMO directly produces the reference coordinates. All 100 targets depart and arrive, and their realized edge sequences match their prescribed plans. Predetermined continuity and stop checks accept 97 references: 22 buses, 18 motorcycles and 19 in each remaining class. Two bus requests fail, one for a stop-count mismatch and one for a coordinate jump; one motorcycle fails the continuity or validity check. None is replaced with an easier request.

The generation-side map constructor independently supports 97 of the original 100 requests. Its three failures and the three reference failures are distinct, leaving 94 requests with both a supported map condition and an accepted reference. Thus 97/100 describes generation geometry coverage, 97/100 describes reference acceptance, and 94/100 is the paired quality population. The generator is run without consulting reference acceptance: eight candidates per supported request produce 776 candidates per full-pool model or baseline, of which 752 belong to the 94 paired requests. Selection later yields one trajectory per generated request while retaining the same coverage accounting.

The event record also limits interpretation. The new simulation contains 24 background-and-target teleport starts in total. The one directly affected target bus is already outside the accepted reference cohort under the frozen extraction checks. An accepted ride-hailing target has an emergency-stop warning and remains included. No post-hoc resimulation or removal is used to improve the reference distribution. The reference population is therefore a documented simulator realization with retained limitations, not a claim that every accepted movement is free of simulator anomalies.

The model receives only the frozen map request and a duration sampled from training-side statistics. Reference coordinates, realized duration and future lane choices become available to the evaluator after generation. Their role is to characterize the distribution of outcomes for comparable conditions; they do not identify a unique correct generated sequence. Physical distances are consequently reported as class-wise distribution comparisons with equal request weight, accompanied by support counts. When a class lacks support, the fixed five-class macro is undefined rather than silently recomputed over a smaller set.

Finally, independence of this new simulation has a limited meaning. It has a new stochastic realization, target identities and departure instances, but shares the map and public demand families with training. Bus services in particular are familiar templates, not held-out route families. The experiment is a same-map prospective development evaluation on one new episode. It supports controlled comparisons of the generator and GAF feedback under those conditions; it does not supply independent city-level replication or justify confidence intervals that treat every candidate as a separate traffic simulation.
